\documentclass[conference]{IEEEtran}

\usepackage{cite}
\usepackage{amsmath,amssymb}
\usepackage{graphicx}
\usepackage{hyperref}
\usepackage{booktabs}

% ── Metadata ─────────────────────────────────────────────────────────────────
\title{From Interpretable to Black-Box: Dynamic Model Identification for Small USV Control Design}

\author{
  \IEEEauthorblockN{Brian S. Bingham}
  \IEEEauthorblockA{
    Department of Mechanical and Aerospace Engineering\\
    Naval Postgraduate School\\
    Monterey, CA, USA\\
    bbingham@nps.edu
  }
}

\begin{document}

\maketitle

% ── Abstract ─────────────────────────────────────────────────────────────────
\begin{abstract}
Effective low-level feedback control of small underactuated uncrewed surface
vehicles (USVs), specifically vessels under seven meters with a single thruster
and rudder, requires dynamic models that provide both actionable design insight and 
predictive accuracy. Tow-tank testing is impractical at this scale;
parameters must be identified solely from limited field trials. Existing system
identification literature largely pursues model fidelity as an end in itself,
producing representations whose complexity can obscure which aspects of the
dynamics are structurally significant for compensator design.

We propose a methodology built around a deliberate spectrum of three model levels,
each identified from the same experimental data: (1) first-order, linear,
single-input single-output (SISO) transfer functions, Nomoto-style models for
surge and yaw that are immediately interpretable and support controller architecture
intuition, but risk oversimplification when considered in isolation;
(2) a physics-based nonlinear state-space formulation incorporating quadratic
hydrodynamic drag; and (3) a Nonlinear AutoRegressive with eXogenous inputs (NARX)
multilayer perceptron. Rather than seeking a single best model, the spectrum is
used collectively: simple models anchor human intuition; structured nonlinear models
expose which terms are dynamically salient; and the black-box model bounds the
benefit of added complexity.

Preliminary results from step-response field experiments demonstrate the approach.
First-order SISO models achieve R2=0.93 (surge, tau=1.95 s) and R2=0.83
(yaw rate, tau=0.47 s), providing direct controller tuning parameters.
Cross-model diagnostics identify quadratic drag as the dominant surge nonlinearity
and reveal that speed-dependent rudder effectiveness, a key coupling for
underactuated control design, is not resolvable from a single short trial,
motivating targeted data collection. The NARX model matches one-step prediction
accuracy (R2=0.93) but diverges in closed-loop simulation, demonstrating that
physical structure provides essential regularization under limited data. A
comprehensive dataset spanning multiple operating regimes is being collected to
evaluate the full methodology.
\end{abstract}

\begin{IEEEkeywords}
uncrewed surface vehicle, system identification, model hierarchy, feedback control
design, nonlinear dynamics, data-driven modeling, underactuated vessels
\end{IEEEkeywords}

% ─────────────────────────────────────────────────────────────────────────────
\section{Introduction}
\label{sec:intro}
% OUTLINE:
% - Motivate the problem: small USVs (<7 m, single thruster + rudder) are
%   increasingly deployed for survey, inspection, and autonomy research.
%   Low-level speed and heading control is a prerequisite for any higher-level
%   autonomy.
%
% - State the tension: control design requires a model, but the "right" model
%   depends on what the model is for. The field tends to treat system ID as
%   pure estimation (maximize fidelity) rather than as a tool for design insight.
%
% - Introduce the Box framing (implicit): a spectrum of models, each wrong in
%   different ways, is collectively more useful than one complex model.
%
% - Scope the paper: single thruster + rudder, <7 m, field-identified parameters.
%   Distinguish from differential-thrust vessels and from large vessels with
%   tow-tank access.
%
% - State contributions:
%   1. A three-level model spectrum methodology for this class of USV
%   2. Cross-model diagnostics that identify structurally salient dynamics
%   3. An empirical demonstration showing when physical structure matters
%      (black-box simulation failure under limited data)
%   4. A targeted data collection framework motivated by diagnostic gaps


% ─────────────────────────────────────────────────────────────────────────────
\section{Related Work}
\label{sec:related}
% OUTLINE:
% - Physics-based USV system identification
%   * Sonnenburg & Woolsey (2013) [journal] and (2010) [OCEANS]: closest prior
%     work -- explicitly compares model hierarchy from 3-DOF nonlinear down to
%     Nomoto first-order; finds first-order best compromise at higher speeds.
%     Does not include data-driven models.
%   * Wang et al. (2020): 3-DOF nonlinear model from sea trials.
%   * Nomoto (1957): foundational first-order steering model.
%   * Fossen (2011): standard reference for USV dynamics formulation.
%
% - Data-driven and ML approaches
%   * Woo et al. (2018): LSTM-RNN outperforms linear models (76.9% surge error
%     reduction); evaluates one-step prediction only, not closed-loop simulation.
%   * Xu et al. (2022): PINN for 3-DOF USV; better generalization with small
%     datasets by embedding physical structure in the loss.
%   * Tanveer & Ahmad (2023): linear second-order yaw model from pool trials.
%
% - Gap: no prior work explicitly constructs a spectrum across all three levels
%   and evaluates it as a collective design tool. Prior work either pursues
%   maximum fidelity or demonstrates a single ML approach, without cross-model
%   diagnostics or honest closed-loop simulation comparison.
%
% - Note on scope: most prior work uses differential-thrust vessels or larger
%   platforms. Single-thruster, rudder-steered USVs are underrepresented and
%   present a harder underactuation problem.


% ─────────────────────────────────────────────────────────────────────────────
\section{Model Spectrum Methodology}
\label{sec:methodology}
% OUTLINE:
% - Framing: three levels, each identified from the same dataset.
%   Each level answers a different design question.

\subsection{Level~1: First-Order Linear SISO Models}
\label{sec:l1}
% OUTLINE:
% - Nomoto-style first-order transfer functions for surge and yaw, treated as
%   independent SISO systems.
%
%   Surge: G_v(s) = K / (tau*s + 1),  input=throttle (normalized), output=speed
%   Yaw:   G_r(s) = K_r / (tau_r*s + 1), input=rudder, output=yaw rate
%
% - Identification: discrete-time ARX(1,1) via ordinary least squares; convert
%   poles to physical parameters tau and K.
%
% - Design utility: directly readable time constants and gains for PID tuning;
%   immediately reveals dominant timescales for each channel.
%
% - Limitation: linear, decoupled; ignores quadratic drag, surge-yaw coupling,
%   and operating-regime dependence. Taken alone, risks oversimplification that
%   leads to a compensator architecture mismatch.

\subsection{Level~2: Nonlinear State-Space Model}
\label{sec:l2}
% OUTLINE:
% - Physics-based, two-state model derived from planar rigid-body equations:
%
%   dv/dt = k_T * u - d2 * v|v|       (surge)
%   dr/dt = k_delta * delta - d_r * r  (yaw)
%
% - Quadratic drag term v|v| is the key nonlinearity relative to Level 1.
%
% - Two-phase identification strategy:
%   Phase 1 (coasting, u~0): identify d2 from zero-thrust deceleration segments.
%   Phase 2 (powered): identify k_T with d2 fixed.
%   This decoupling avoids conflation of thrust and drag in a joint fit.
%
% - Design utility: reveals the drag structure relevant to controller design
%   (e.g., feedforward compensation, speed-hold steady state).
%   Maintains physical interpretability while adding the dominant nonlinearity.

\subsection{Level~3: Black-Box NARX Model}
\label{sec:l3}
% OUTLINE:
% - Nonlinear AutoRegressive with eXogenous inputs (NARX) multilayer perceptron.
%   Predicts [v(k+1), r(k+1)] from n_lags past samples of all four signals:
%   [v, r, throttle, rudder]. MIMO: discovers coupling the structured models
%   assume away.
%
% - Architecture: 3 lags, 1 hidden layer (20 neurons), ReLU activation,
%   StandardScaler normalization, early stopping to guard against overfitting.
%
% - Training: chronological 80/20 split (first 100 s train, last 30 s test).
%   Chronological split avoids temporal leakage.
%
% - Evaluation: one-step prediction AND closed-loop simulation separately
%   reported -- these diverge significantly with limited data.
%
% - Design utility: bounds the benefit of added complexity given available data;
%   motivates richer data collection; a reference point for future comparison
%   as dataset grows.

\subsection{Cross-Model Diagnostics}
\label{sec:diagnostics}
% OUTLINE:
% - The spectrum's value is not just the three fits -- it is the comparative
%   analysis across levels that reveals salient structure.
%
% - Three diagnostics performed:
%   1. Speed-dependent rudder gain: does replacing constant K_r with K_r*v
%      improve yaw-rate fit? (Result: no improvement with current data --
%      ΔR² = -0.03 -- flags gap in data coverage, not absence of effect)
%
%   2. Drag nonlinearity: linear vs. quadratic fit to coasting deceleration.
%      (Result: quadratic dominant, ΔR² = +0.20; D1≈0, D2=0.44 m^{-1})
%
%   3. Rudder-induced surge drag: correlation of δ² with surge residuals.
%      (Result: negligible, r = -0.08; coupling term not warranted)
%
% - Implication: diagnostics directly inform the Level-2 model structure AND
%   identify what the next data collection must target.


% ─────────────────────────────────────────────────────────────────────────────
\section{Experimental Platform and Data Collection}
\label{sec:experiment}
% OUTLINE:
% - Vehicle description: [USV model/name], approx. [length] m, single
%   propeller + rudder, ArduPilot autopilot, GPS + IMU.
%   Inputs: RC throttle (RCOU_C3) and rudder (RCOU_C1) commands, 1100-1900 us PWM.
%   Outputs: GPS speed (5 Hz), IMU yaw rate (137 Hz).
%
% - Data cleaning: multi-rate streams aligned to 10 Hz common grid via
%   linear interpolation; RC commands normalized to [-1,1].
%   One 29.3 s gap (vessel repositioning) handled by segmentation.
%
% - Current dataset: single step-response trial, 130 s working window
%   (150-280 s of a 341 s recording), throttle range 0-39%, speed 0-2.9 m/s.
%   This dataset is intentionally minimal -- a proof-of-concept for the workflow.
%
% - Planned dataset (full paper): multiple trials spanning throttle 0-100%,
%   dedicated rudder sweeps at fixed speeds (to resolve speed-dependent gain),
%   and maneuvers spanning displacement, semi-displacement, and planing regimes.
%   [describe vehicle operating speed range and transition speeds]


% ─────────────────────────────────────────────────────────────────────────────
\section{Results}
\label{sec:results}

\subsection{Level~1: First-Order SISO}
% OUTLINE:
% - Table: identified parameters (tau, K) for surge and yaw.
% - Simulation R² for both channels.
% - Figure: measured vs. simulated speed and yaw rate.
% - Note: parameters directly usable for PID design.

\subsection{Level~2: Nonlinear State-Space}
% OUTLINE:
% - Table: identified parameters (k_T, d2, k_delta, d_r).
% - Two-phase identification: compare d2 from coasting vs. joint fit;
%   motivate why two-phase is necessary.
% - Simulation R² comparison vs. Level 1.
% - Figure: measured vs. simulated (combined plot with all three levels).

\subsection{Level~3: NARX}
% OUTLINE:
% - One-step prediction R² vs. closed-loop simulation R² -- highlight the gap.
% - Figure: simulation divergence for surge.
% - Interpretation: limited data makes black-box models unreliable simulators;
%   physical structure is essential regularization at this data scale.

\subsection{Cross-Model Diagnostic Summary}
% OUTLINE:
% - Summary table: three diagnostics, findings, implication for model/control design.
% - Key message: the spectrum reveals structure that no single model exposes.


% ─────────────────────────────────────────────────────────────────────────────
\section{Discussion}
\label{sec:discussion}
% OUTLINE:
% - What the spectrum reveals for this class of vehicle:
%   * Surge and yaw are weakly coupled at the operating speeds tested --
%     decoupled SISO controllers are a reasonable starting point.
%   * Quadratic drag is significant -- speed-hold controllers need nonlinear
%     feedforward; linear controllers will have speed-dependent steady-state error.
%   * Speed-dependent rudder gain cannot yet be confirmed or denied --
%     this is the primary gap the expanded dataset must fill.
%     If confirmed, it implies gain scheduling or adaptive control is warranted.
%
% - Implications for compensator architecture selection:
%   * Level 1 → baseline PID structure with identified gains.
%   * Level 2 → nonlinear feedforward for drag compensation;
%     potentially gain-scheduled rudder loop.
%   * Level 3 → motivates MPC or learning-based control once data is sufficient.
%
% - On fidelity vs. insight: even the highest-fidelity model here cannot
%   eliminate experimental fine-tuning. The spectrum sets realistic expectations
%   and guides where to spend tuning effort.
%
% - Limitations of the current dataset and what the full study will add.


% ─────────────────────────────────────────────────────────────────────────────
\section{Conclusion}
\label{sec:conclusion}
% OUTLINE:
% - Summarize the three-level spectrum methodology for small underactuated USVs.
% - Restate the core finding: the spectrum collectively provides more control
%   design guidance than any single model; simple models are not merely a
%   stepping stone but a persistent design tool.
% - State what the full dataset will allow: resolution of speed-dependent
%   rudder dynamics, validation across operating regimes, and a more
%   definitive comparison of black-box vs. structured models.
% - Call to action: release the dataset and analysis notebooks as open tools
%   for the USV research community.


% ─────────────────────────────────────────────────────────────────────────────
\nocite{*}  % include all references even while body is outline-only
\bibliographystyle{IEEEtran}
\bibliography{references}

\end{document}
